\documentclass[10pt,a4paper]{amsart}
\usepackage[margin=19mm]{geometry}
\usepackage{amsmath,amssymb,mathtools}
\usepackage[scr=rsfs,cal=euler]{mathalfa}
\usepackage{xfrac}
\usepackage[unicode,hidelinks,bookmarks=false]{hyperref}
\newcommand{\ie}{i.e.\ }
\newcommand{\cf}{cf.\ }
\newcommand{\resp}{respectively\ }
\newcommand{\Subsection}{\subsection}
\providecommand{\nobreakdash}{\nobreak\hbox{-}}
\setlength{\emergencystretch}{2em}
\multlinegap0pt
\usepackage{natbib}
\usepackage{multirow}
\usepackage{amssymb}
\usepackage{xcolor}
\usepackage{booktabs}
\usepackage{algorithm}\usepackage{algpseudocode}
\usepackage{caption}
\newtheorem{theorem}{Theorem}
\title{On the maximum $\sigma$-irregularity of trees with given order and maximum degree}
\author{Milan Ba\v{s}i\'c}
\date{Source: arXiv:2602.14241v1 (CC BY 4.0)}
\begin{document}
\maketitle

\noindent\emph{Source and licence.} On the maximum $\sigma$-irregularity of trees with given order and maximum degree. Version arXiv:2602.14241v1 (CC BY 4.0), distributed by arXiv under the Creative Commons Attribution 4.0 International licence, http://creativecommons.org/licenses/by/4.0/, as recorded in the arXiv licence metadata for this exact version and shown on the abstract page https://arxiv.org/abs/2602.14241v1. Full licence notice: \texttt{licenses/arxiv-2602.14241-CC-BY-4.0.txt}.\par\smallskip

\noindent\textit{Editorial scope.} Counted unit: the article's theorem thm:SO-LP-1-Delta (source0.tex lines 424--451). The article's complete original proof of it is delivered with it (source0.tex lines 453--734, one \texttt{proof} environment of 282 lines). No mathematics is written, repaired or re-derived: the statement and the proof are the article's own lines, in the article's own notation, and no retained line is substituted or re-flowed.  No further source-local context is needed: every cross-reference of every retained line resolves inside the two delivered spans.  Nothing was omitted from any retained span, so the package carries no omission record.  No retained line contains a citation, so the wrapper carries no bibliography and no bibliography entry.  No retained line contains a figure environment, an image-inclusion command or drawing code, so no image is delivered and the package declares no asset dependency.  Editorial formatting only, in addition: an AMS \texttt{amsart} preamble that restates the article's own declarations and macros used by the retained lines, namely the environments \texttt{theorem} on separate counters, together with the standard packages those lines need.  The retained environments print in this document as Theorem 1 in order of appearance, whereas the article prints the same items with the numbering of its own full document.  The complete native source of the article as distributed in the arXiv e-print for this exact version is bundled unchanged as \texttt{sources/194698/source0.tex}.
\begin{theorem}
\label{thm:SO-LP-1-Delta}
Let $\Delta\ge 4$ and $n\ge \Delta+2$ be integers. Consider the linear program
\begin{align}
\max\ & \sigma(T) \;=\; \sum_{1\le i\le j\le \Delta} (i-j)^2\, m_{i,j},
\label{eq:primal-objective}\\[1ex]
\text{subject to }&
\sum_{1\le i\le j\le \Delta} \Bigl(\frac{1}{i}+\frac{1}{j}\Bigr) m_{i,j} \;=\; n,
\label{eq:primal-constraint-w}\\
&
\sum_{1\le i\le j\le \Delta} m_{i,j} \;=\; n-1,
\label{eq:primal-constraint-sum}\\
&
m_{i,j}\;\ge\; 0
\quad\text{for all }1\le i\le j\le \Delta.
\label{eq:primal-nonneg}
\end{align}
Then the problem admits a unique optimal solution in which
\[
m_{1,\Delta} \;=\; \frac{(\Delta-2)n + (\Delta+2)}{\Delta}, 
\qquad
m_{2,\Delta} \;=\; \frac{2(n-\Delta-1)}{\Delta},
\]
and
\[
m_{i,j}=0 \quad\text{for all } (i,j)\notin\{(1,\Delta),(2,\Delta)\}.
\]
\end{theorem}

\begin{proof}
Define
\[
w_{i,j} := \frac{1}{i}+\frac{1}{j},
\qquad
\sigma_{i,j} := (i-j)^2.
\]
Then the primal problem \eqref{eq:primal-objective}--\eqref{eq:primal-nonneg}
can be written in the standard form
\[
\max \sum_{1\le i\le j\le \Delta} \sigma_{i,j} m_{i,j}
\quad\text{subject to}\quad
\sum_{i\le j} w_{i,j} m_{i,j} = n,\;
\sum_{i\le j} m_{i,j} = n-1,\;
m_{i,j}\ge 0.
\]

\medskip\noindent
\emph{Dual problem.}
The dual linear program associated with this primal problem is
\begin{equation}
\label{eq:dual-problem}
\begin{aligned}
\min\ & D(\lambda,\mu) \;=\; n\lambda + (n-1)\mu,\\
\text{subject to }&
\lambda w_{i,j} + \mu \;\ge\; \sigma_{i,j},
\quad 1\le i\le j\le \Delta,
\end{aligned}
\end{equation}
where $\lambda,\mu\in\mathbb{R}$ are the dual variables corresponding to
constraints \eqref{eq:primal-constraint-w} and \eqref{eq:primal-constraint-sum},
respectively.

For convenience, we define the dual slack function
\begin{equation}
\label{eq:Fij-def}
F(i,j)
\;:=\;
\lambda w_{i,j} + \mu - \sigma_{i,j}.
\end{equation}
Dual feasibility is equivalent to $F(i,j)\ge 0$ for all $1\le i\le j\le \Delta$.

\medskip\noindent
\emph{Choice of dual variables.}
We now determine $\lambda$ and $\mu$ by requiring that the dual constraints
be tight for the index pairs $(1,\Delta)$ and $(2,\Delta)$, that is,
\begin{equation}
\label{eq:dual-equalities-1-2}
F(1,\Delta)=0,
\qquad
F(2,\Delta)=0.
\end{equation}
Equivalently,
\[
\lambda w_{1,\Delta} + \mu = \sigma_{1,\Delta} = (\Delta-1)^2,
\qquad
\lambda w_{2,\Delta} + \mu = \sigma_{2,\Delta} = (\Delta-2)^2.
\]

Since
\[
w_{1,\Delta} = 1+\frac{1}{\Delta},
\qquad
w_{2,\Delta} = \frac12+\frac{1}{\Delta},
\]
subtracting the two equalities gives
\[
(\Delta-1)^2 - (\Delta-2)^2
=
\lambda\Bigl(1+\frac{1}{\Delta}-\frac12-\frac{1}{\Delta}\Bigr)
=
\frac{\lambda}{2},
\]
hence
\begin{equation}
\label{eq:lambda}
\lambda
=
2\bigl((\Delta-1)^2-(\Delta-2)^2\bigr)
=
4\Delta-6.
\end{equation}
Substituting this value into the first equality yields
\begin{equation}
\label{eq:mu}
\mu
=
(\Delta-1)^2 - (4\Delta-6)\Bigl(1+\frac{1}{\Delta}\Bigr)
=
\Delta^2 - 6\Delta + 3 + \frac{6}{\Delta}.
\end{equation}

With this choice of $\lambda$ and $\mu$, the slack function
\eqref{eq:Fij-def} takes the explicit form
\begin{equation}
\label{eq:Fij-explicit}
F(i,j)
=
(4\Delta-6)\Bigl(\frac{1}{i}+\frac{1}{j}\Bigr)
+ \Delta^2 - 6\Delta + 3 + \frac{6}{\Delta}
- (i-j)^2.
\end{equation}

\medskip\noindent
\emph{Verification of dual feasibility.}

\smallskip\noindent
\emph{Case 1: $i=1$, $1\le j\le \Delta$.}
From \eqref{eq:Fij-explicit},
\[
F(1,j)
=
(4\Delta-6)\Bigl(1+\frac{1}{j}\Bigr)
+ \Delta^2 - 6\Delta + 3 + \frac{6}{\Delta}
- (1-j)^2.
\]
A direct computation yields
\begin{equation}
\label{eq:F1j-factor}
\Delta j\,F(1,j)
=
(\Delta-j)
\left(
\Delta j^2 + \Delta^2 j - 2\Delta j + 4\Delta - 6
\right).
\end{equation}
For $1\le j\le \Delta-1$ we have $\Delta-j>0$, so it suffices to analyze
\[
G_1(j)
=
\Delta j^2 + (\Delta^2-2\Delta)j + 4\Delta - 6.
\]
This quadratic polynomial is strictly convex, with vertex at
\[
j_0=-\frac{\Delta^2-2\Delta}{2\Delta}=-\frac{\Delta-2}{2}\le 0,
\]
and hence is strictly increasing on $[0,\infty)$. Therefore,
\[
\min_{1\le j\le \Delta-1} G_1(j)=G_1(1)=\Delta^2+3\Delta-6>0
\quad(\Delta\ge 4).
\]
It follows from \eqref{eq:F1j-factor} that
\[
F(1,j)>0 \quad\text{for }1\le j\le \Delta-1,
\qquad
F(1,\Delta)=0.
\]

\smallskip\noindent
\emph{Case 2: $i=2$, $2\le j\le \Delta$.}
Similarly,
\[
F(2,j)
=
(4\Delta-6)\Bigl(\frac12+\frac{1}{j}\Bigr)
+ \Delta^2 - 6\Delta + 3 + \frac{6}{\Delta}
- (2-j)^2,
\]
and a direct computation yields
\begin{equation}
\label{eq:F2j-factor}
\Delta j\,F(2,j)
=
(\Delta-j)
\left(
\Delta j^2 + \Delta^2 j - 4\Delta j + 4\Delta - 6
\right).
\end{equation}
For $2\le j\le \Delta-1$, consider
\[
G_2(j)
=
\Delta j^2 + (\Delta^2-4\Delta)j + 4\Delta - 6.
\]
This is again strictly convex with vertex
\[
j_0=-\frac{\Delta^2-4\Delta}{2\Delta}=-\frac{\Delta-4}{2}\le 0,
\]
so $G_2$ is increasing on $[0,\infty)$ and
\[
\min_{2\le j\le \Delta-1} G_2(j)=G_2(2)=2\Delta^2-6>0.
\]
Hence
\[
F(2,j)>0 \quad\text{for }2\le j\le \Delta-1,
\qquad
F(2,\Delta)=0.
\]

\smallskip\noindent
\emph{Case 3: $3\le i\le j\le \Delta$.}
Using $1/i\ge 1/\Delta$ and $1/j\ge 1/\Delta$, we obtain
\[
F(i,j)
\ge
(4\Delta-6)\frac{2}{\Delta}
+ \Delta^2 - 6\Delta + 3 + \frac{6}{\Delta}
- (i-j)^2.
\]
Since $0\le j-i\le \Delta-3$, we have $(i-j)^2\le (\Delta-3)^2$, and therefore
\begin{align*}
F(i,j)
&\ge
\Delta^2 - 6\Delta + 3 + \frac{6}{\Delta}
+ \frac{8\Delta-12}{\Delta}
- (\Delta-3)^2\\
&=
2-\frac{6}{\Delta}>0
\qquad(\Delta\ge 4).
\end{align*}

Combining the three cases, we conclude that
\[
F(i,j)\ge 0 \quad\text{for all }1\le i\le j\le \Delta,
\]
with equality if and only if $(i,j)\in\{(1,\Delta),(2,\Delta)\}$.
Thus $(\lambda,\mu)$ is dual feasible.

\medskip\noindent
\emph{Construction of a primal feasible solution.}
Assume $m_{i,j}=0$ for all $(i,j)\notin\{(1,\Delta),(2,\Delta)\}$.
Then constraints \eqref{eq:primal-constraint-w}--\eqref{eq:primal-constraint-sum}
reduce to
\begin{equation}
\label{eq: reduced system}
m_{1,\Delta}+m_{2,\Delta}=n-1,
\qquad
w_{1,\Delta}m_{1,\Delta}+w_{2,\Delta}m_{2,\Delta}=n.
\end{equation}
Solving this system yields
\[
m_{1,\Delta}=\frac{(\Delta-2)n+(\Delta+2)}{\Delta},
\qquad
m_{2,\Delta}=\frac{2(n-\Delta-1)}{\Delta}.
\]
Since $n\ge\Delta+2$, both values are positive.

\medskip\noindent
\emph{Optimality and uniqueness.}
For the primal solution constructed above, we compute
\begin{align*}
\sigma(T)
&=
\sum_{1\le i\le j\le \Delta} \sigma_{i,j} m_{i,j}\\
&=
\sigma_{1,\Delta}m_{1,\Delta}+\sigma_{2,\Delta}m_{2,\Delta}
\qquad(\text{all other }m_{i,j}=0)\\
&=
(\lambda w_{1,\Delta}+\mu)m_{1,\Delta}+(\lambda w_{2,\Delta}+\mu)m_{2,\Delta}
\qquad(\text{by }F(1,\Delta)=F(2,\Delta)=0)\\
&=
\sum_{1\le i\le j\le \Delta} (\lambda w_{i,j}+\mu)\, m_{i,j}\\
&=
\lambda\sum_{i\le j} w_{i,j} m_{i,j}
+
\mu\sum_{i\le j} m_{i,j}
=
\lambda n + \mu (n-1)
=
D(\lambda,\mu),
\end{align*}
which coincides with the dual objective value at the feasible dual solution
$(\lambda,\mu)$ given in \eqref{eq:lambda} and \eqref{eq:mu}.
By strong duality, both the primal and dual solutions are optimal.

Furthermore, we have shown that
\[
F(i,j)=\lambda w_{i,j}+\mu-\sigma_{i,j}>0
\quad\text{for all }(i,j)\notin\{(1,\Delta),(2,\Delta)\},
\]
while $F(1,\Delta)=F(2,\Delta)=0$.
By complementary slackness, this implies that in any optimal primal solution
one must have
\[
m_{i,j}=0
\quad\text{for all }(i,j)\notin\{(1,\Delta),(2,\Delta)\}.
\]
Since the reduced system \eqref{eq: reduced system}
has a unique solution, the optimal primal solution is unique.
This completes the proof.

\end{proof}

\end{document}
